% -*- coding: utf-8 -*-
% !TEX program = xelatex
\documentclass[plain,noanswer,medmath]{../../template/jnuexam}
\usepackage{../../template/kysx}

\begin{document}


\examtitle{1997}{4}


\loadproblems{3}{1997P3}%载入试卷三


\makepart{填空题}{本题共5分，每小题3分，满分15分}


\useproblem{3}{1}{1}%【同试卷三第一(1)题】


\begin{problem}
设$f(x) = \frac{1}{1 + x^2} + x^3\int_0^1f(x)\dx$，则$\int_0^1f(x)\dx =$ \fillin{}．
\end{problem}


\begin{problem}
设$n$阶矩阵$A = \begin{pmatrix}
       0 &      1 &      1 & \cdots &      1 & 1 \\
       1 &      0 &      1 & \cdots &      1 & 1 \\
       1 &      1 &      0 & \cdots &      1 & 1 \\
  \vdots & \vdots & \vdots & \ddots & \vdots & \vdots \\
       1 &      1 &      1 & \cdots &      0 & 1 \\
       1 &      1 &      1 & \cdots &      1 & 0
\end{pmatrix}$，则$|A|=$ \fillin{}．
\end{problem}


\begin{problem}
设$A$, $B$是任意两个随机事件，则$P\{(\bar A + B)(A + B)(\bar A + \bar B)(A + \bar B)\} =$ \fillin{}．
\end{problem}


\begin{problem}
设随机变量$X$服从参数为$(2,p)$的二项分布，随机变量$Y$服从参数为$(3,p)$的二项分布．
若$P\{X \ge 1\} = \frac{5}{9}$，则$P\{Y \ge 1\} =$ \fillin{}．
\end{problem}


\makepart{选择题}{本题共5小题，每小题3分，满分15分}


\begin{problem}
设$f(x)$, $\phi (x)$在点$x = 0$的某邻域内连续，且当$x \to 0$时，$f(x)$是$\phi (x)$的高阶无穷小，
则当$x \to 0$时，$\int_0^x f(t)\sin t\dt$是$\int_0^x t\phi(t)\dt$的\pickout{}
\begin{abcd}
\item 低阶无穷小．                    
\item 高阶无穷小．
\item 同阶但不等价的无穷小．          
\item 等价无穷小．
\end{abcd}
\end{problem}


\useproblem{3}{2}{2}%【同试卷三第二(2)题】


\useproblem{3}{2}{3}%【同试卷三第二(3)题】


\begin{problem}
非齐次线性方程组$AX = b$中未知量个数为$n$，方程个数为$m$，系数矩阵$A$的秩为$r$，则\pickout{}
\begin{abcd}
\item $r = m$时，方程组$AX = b$有解．
\item $r = n$时，方程组$AX = b$有唯一解．
\item $m = n$时，方程组$AX = b$有唯一解．
\item $r < n$时，方程组$AX = b$有无穷多解．
\end{abcd}
\end{problem}


\begin{problem}
设$X$是一随机变量，$EX = \mu$, $DX = \sigma^2$（$\mu$, $\sigma > 0$是常数），则对任意常数$c$，必有\pickout{}
\begin{abcd}
\item $E(X - c)^2 = E(X^2) - c^2$．			
\item $E(X - c)^2 = E(X - \mu)^2$．
\item $E(X - c)^2 < E(X - \mu)^2$．			
\item $E(X - c)^2 \ge E(X - \mu)^2$．
\end{abcd}
\end{problem}


\makepart{}{本题满分6分}

\begin{problem}
求极限$\lim_{x \to 0} \left[\frac{a}{x} - \left(\frac{1}{x^2} - a^2 \right)\ln(1 + ax) \right]$%
（$a \ne 0$）．
\end{problem}


\makepart{}{本题满分6分}%【同试卷三第四题】

\useproblem{3}{4}{0}


\makepart{}{本题满分6分}

\begin{problem}
假设某种商品的需求量$Q$是单价$p$（单位:元）的函数：$Q = 12000 - 80p$；
商品的总成本$C$是需求量$Q$的函数：$C = 25000 + 50Q$；
每单位商品需要纳税 $2$元．试求使销售利润最大的商品单价和最大利润额．
\end{problem}


\makepart{}{本题满分7分}

\begin{problem}
求曲线$y = x^2 - 2x$, $y = 0$, $x = 1$, $x = 3$所围成的平面图形的面积$S$，
并求该平面图形绕$y$轴旋转一周所得旋转体的体积$V$．
\end{problem}


\makepart{}{本题满分7分}

\begin{problem}
设函数$f(x)$在$(-\infty , +\infty)$内连续，且$F(x) = \int_0^x (x - 2t) f(t)\dt$．试证：
\begin{enumerate}
\item 若$f(x)$为偶函数，则$F(x)$也是偶函数；
\item 若$f(x)$单调不增，则$F(x)$单调不减．
\end{enumerate}
\end{problem}


\makepart{}{本题满分6分}

\begin{problem}
设$D$是以点$O(0,0)$, $A(1,2)$和$B(2,1)$为顶点的三角形区域，求$\iint_Dx\dx\dy$．
\end{problem}


\makepart{}{本题满分7分}%【同试卷三第九题】

\useproblem{3}{9}{0}


\makepart{}{本题满分9分}

\begin{problem}
设矩阵$A$和$B$相似，且$A = \begin{pmatrix}
    1 & - 1 & 1 \\
    2 &   4 & - 2 \\
  - 3 & - 3 & a
\end{pmatrix}$，$B = \begin{pmatrix}
  2 & 0 & 0 \\
  0 & 2 & 0 \\
  0 & 0 & b
\end{pmatrix}$，\par
\begin{enumerate*}
\item 求$a$, $b$的值；
\item 求可逆矩阵$P$，使$P^{-1} AP = B$．
\end{enumerate*}
\end{problem}


\makepart{}{本题满分8分}%【类似试卷三第十一题】

\begin{problem}
假设随机变量$X$的绝对值不大于$1$；$P\{X = - 1\} = \frac{1}{8},P\{X = 1\} = \frac{1}{4}$；
在事件$\{- 1 < X < 1\}$出现的条件下，$X$在$(- 1,1)$内的任一子区间上取值的条件概率与该子区间长度成正比．
试求：\par
\begin{enumerate*}
\item $X$的分布函数$F(x) = P\{X \le x\}$；
\item $X$取负值的概率$p$．
\end{enumerate*}
\end{problem}


\makepart{}{本题满分8分}

\begin{problem}
假设随机变量$Y$服从参数为$\lambda = 1$的指数分布，随机变量$X_k = \begin{cases}
  0, & Y \le k, \\
  1, & Y > k
\end{cases}$（$k = 1,2$），求：
\begin{enumerate*}
\item $X_1$和$X_2$的联合概率分布；
\item $E(X_1 + X_2)$．
\end{enumerate*}
\end{problem}


\end{document}
